\documentclass[11pt,a4paper]{article}
\usepackage{fontspec}
\setmainfont{DejaVu Sans}[Scale=0.90]
\usepackage{amsmath,amssymb}
\usepackage[margin=2.4cm]{geometry}
\usepackage{booktabs,array,longtable,tabularx}
\usepackage{enumitem}
\usepackage{titlesec}
\usepackage{parskip}
\titlelabel{\thetitle.\quad}
\titleformat{\section}{\normalfont\large\bfseries}{\thesection.}{0.6em}{}
\titleformat{\subsection}{\normalfont\normalsize\bfseries}{\thesubsection}{0.6em}{}
\titlespacing*{\section}{0pt}{16pt}{6pt}
\titlespacing*{\subsection}{0pt}{12pt}{4pt}
\setlist[itemize]{leftmargin=1.4em,itemsep=1pt,topsep=3pt}
\newcolumntype{L}[1]{>{\raggedright\arraybackslash}p{#1}}
\begin{document}
\begin{center}
{\Large\bfseries Compréhension de l'article --- Krapez (2018)}\\[4pt]
{\small Krapez, \emph{International Journal of Thermal Sciences} \textbf{136} (2018) 182--199}\\[2pt]
{\small Version corrigée --- 8 septembre 2026}
\end{center}

space{4pt}


\section{Les deux équations de départ}

\subsection{Formulation duale}

L'article traite simultanément deux équations. La première porte sur la température transformée,
la seconde sur la densité de flux transformée :
\[ p\,\rho c(z)\,\theta = \frac{d}{dz}\!\left[\lambda(z)\,\frac{d\theta}{dz}\right] ,
\qquad
p\,\frac{1}{\lambda(z)}\,\phi = \frac{d}{dz}\!\left[\frac{1}{\rho c(z)}\,\frac{d\phi}{dz}\right] . \]

Dans la première, la conductivité figure sous l'opérateur de dérivation et la capacité volumique
intervient algébriquement. Dans la seconde, les rôles sont permutés et les propriétés apparaissent
sous forme inverse.

Cette dualité n'introduit aucune physique nouvelle. Sa conséquence est que toute construction menée sur
l'une des deux formulations produit simultanément un second profil de matériau, distinct du premier.
Le catalogue de profils solvables est ainsi doublé sans travail supplémentaire.

\subsection{Origine de la formulation en flux}

La formulation en flux s'obtient en éliminant la température au lieu du flux. Les deux relations de
départ sont la loi de Fourier et le bilan d'énergie :
\[ \varphi = -\lambda(z)\,\frac{\partial T}{\partial z} ,
\qquad
\rho c(z)\,\frac{\partial T}{\partial t} = -\frac{\partial\varphi}{\partial z} . \]

En isolant les deux dérivées de $T$ et en utilisant l'égalité des dérivées croisées,
$\partial_t(\partial_z T) = \partial_z(\partial_t T)$, la température disparaît et il reste
\[ \frac{1}{\lambda(z)}\,\frac{\partial\varphi}{\partial t}
= \frac{\partial}{\partial z}\!\left[\frac{1}{\rho c(z)}\,\frac{\partial\varphi}{\partial z}\right] . \]

L'inversion algébrique du bilan qui rend cette élimination possible est propre à la géométrie
unidimensionnelle. En dimension supérieure, l'opération met en jeu l'inversion d'une divergence.

\subsection{Passage au domaine spectral}

La séparation $T(z,t) = \theta(z)\,e^{pt}$ transforme la dérivée temporelle en une multiplication par
$p$, puisque l'exponentielle est proportionnelle à sa propre dérivée. L'exponentielle se simplifie des
deux membres et l'équation aux dérivées partielles devient une équation différentielle ordinaire en
$z$, paramétrée par $p$.

En régime harmonique, $p = i\omega$ et $\theta$ est l'amplitude complexe de la température : son module
donne l'amplitude d'oscillation, son argument le déphasage. En régime transitoire, $p$ est la variable
de Laplace et $\theta$ est une transformée intégrale, sans interprétation ponctuelle équivalente.

\section{Changement de profondeur}

\subsection{Définition de la coordonnée transformée}

\[ \xi(z) = \int_{0}^{z} \frac{du}{\sqrt{a(u)}} , \qquad [\xi] = \mathrm{s^{1/2}} . \]

L'analyse dimensionnelle donne $[a] = \mathrm{m^{2}s^{-1}}$, donc $[a^{-1/2}] = \mathrm{s^{1/2}m^{-1}}$
et $[\xi] = \mathrm{s^{1/2}}$. Le carré de $\xi$ possède ainsi la dimension d'un temps : le temps de
diffusion cumulé le long du trajet. Pour un milieu homogène, $\xi = z/\sqrt{a}$ et $\xi^{2} = z^{2}/a$.

\subsection{Justification de l'exposant}

L'équation de la chaleur est du second ordre en espace et du premier ordre en temps ; le groupement
caractéristique est donc $z^{2}/a$. L'exposant $-1/2$ est le seul qui absorbe entièrement la
diffusivité dans la graduation de l'axe spatial.

Un exposant $-1$ conduirait à la résistance thermique intégrée, qui définit une autre coordonnée et une
autre structure de solutions, associée aux fonctions de Bessel.

\subsection{Équations transformées}

En appliquant $d/dz = a^{-1/2}\,d/d\xi$ et les identités $\lambda/\sqrt{a} = b$ et
$\rho c = b/\sqrt{a}$, les deux équations deviennent
\[ p\,b\,\theta = \bigl(b\,\theta'\bigr)' , \qquad p\,b^{-1}\phi = \bigl(b^{-1}\phi'\bigr)' , \]
où l'accent désigne désormais la dérivation par rapport à $\xi$.

La diffusivité n'apparaît plus dans ces équations. Elle n'a pas disparu de la physique : elle a été
transférée dans la définition de la coordonnée. Le nombre de fonctions matériau inconnues passe ainsi
de deux à une, l'effusivité. Cette réduction constitue le pivot de la méthode : elle rend possible la
démarche inverse consistant à choisir une forme résoluble et à remonter au profil correspondant.

\section{Changement d'inconnue}

\subsection{Transformation de Liouville}

Le développement de l'équation transformée fait apparaître un terme en dérivée première,
\[ \theta'' + \frac{b'}{b}\,\theta' = p\,\theta . \]

La substitution $\theta = \psi\,u$ suivie de l'annulation du coefficient de $\psi'$ conduit à
l'équation $2u' + (b'/b)u = 0$, dont la solution est $u = b^{-1/2}$ à une constante près. Le facteur
n'est donc pas postulé mais déterminé : aucune autre fonction n'annule ce terme.

\subsection{Forme normale}

\[ \frac{d^{2}\psi}{d\xi^{2}} - \bigl(V(\xi) + p\bigr)\psi = 0 ,
\qquad V(\xi) = \frac{s''(\xi)}{s(\xi)} , \qquad s(\xi) = b^{\pm 1/2}(\xi) . \]

L'exposant positif correspond à la formulation en température, l'exposant négatif à la formulation en
flux. Cette convention de double signe est employée dans l'ensemble de l'article.

La dimension de $V$ est celle de l'inverse d'un temps, ce qu'impose son addition à $p$ dans l'équation.

\subsection{Statut de la fonction auxiliaire}

La fonction $\psi$ n'est pas une grandeur physique mesurée. Elle constitue un intermédiaire de calcul
dont l'intérêt est de placer l'équation sous forme normale, ce qui donne accès aux méthodes de
construction de potentiels solvables, notamment les transformations de Darboux.

\section{Prééminence de l'effusivité}

\subsection{Loi de Fourier transformée}

\[ \phi = -\lambda\,\frac{d\theta}{dz}
= -\frac{\lambda}{\sqrt{a}}\,\frac{d\theta}{d\xi}
= -\,b(\xi)\,\theta' . \]

Dans l'espace transformé, la conductivité est remplacée par l'effusivité. Toute condition aux limites
faisant intervenir un flux conductif ne dépend donc que de $b$.

\subsection{Identité des réponses de surface}

L'équation ne contient que $b(\xi)$ ; les conditions aux limites ne contiennent que $b(\xi)$ ; la
solution $\theta(\xi)$ ne dépend donc que de $b(\xi)$.

La surface correspond à $\xi = 0$ pour tout matériau. Deux milieux différents dans l'espace physique
mais partageant le même profil transformé produisent par conséquent une réponse de surface identique :
même amplitude, même phase, même contraste, même effusivité apparente, à toute fréquence et à tout
instant.

La différence entre ces deux milieux correspond à un étirement local différent de l'axe des
profondeurs, sans manifestation observable depuis la surface.

\subsection{Portée du résultat}

Ce résultat de non-unicité concerne le transfert unidimensionnel. En dimension supérieure, la
conduction latérale réintroduit séparément certaines propriétés et la réduction à la seule effusivité
cesse d'être valable.

\section{Méthode des quadripôles}

\subsection{Vecteur d'état}

Le couple retenu est constitué des deux grandeurs continues à une interface parfaite :
\[ \bigl(\theta,\; \phi\bigr) , \qquad \phi = -\,b\,\theta' . \]

La continuité de la température résulte du contact parfait, celle du flux de la conservation de
l'énergie. Le gradient thermique, lui, n'est pas continu lorsque la conductivité présente un saut ; le
couple $(\theta, d\theta/dz)$ ne conviendrait donc pas à l'enchaînement matriciel.

\subsection{Construction de la matrice}

La solution générale de la forme normale s'écrit comme combinaison de deux solutions indépendantes,
$\psi = c_{1}K + c_{2}P$. Le retour aux grandeurs physiques s'effectue par
\[ \theta = \frac{\psi}{s} , \qquad \phi = -\,W(s,\psi) ,
\qquad W(f,g) = f g' - f' g . \]

Le vecteur d'état dépend linéairement des deux constantes, $\mathbf{v} = \mathbf{A}(\xi)\,\mathbf{c}$.
En écrivant cette relation aux deux bords et en éliminant les constantes,
\[ \mathbf{M} = \mathbf{A}(\xi_{0})\,\mathbf{A}(\xi_{1})^{-1} . \]

\subsection{Composition et déterminant}

La relation de transfert étant linéaire et le vecteur de sortie d'une couche étant le vecteur d'entrée
de la suivante, la composition d'un empilement correspond au produit ordonné des matrices,
\[ \mathbf{M}_{\text{total}} = \mathbf{M}_{1}\mathbf{M}_{2}\cdots\mathbf{M}_{N} , \]
l'ordre suivant le parcours de la face avant vers la face arrière.

Pour une équation en forme normale, le théorème d'Abel garantit que le wronskien de deux solutions est
constant. Il en résulte $\det\mathbf{M} = 1$.

\subsection{Limite homogène}

Lorsque l'effusivité devient constante, $V = 0$ et l'équation se réduit à $\psi'' = p\psi$. Les entrées
du quadripôle prennent alors la forme classique
\[ A = D = \cosh\bigl(\sqrt{p}\,\xi_{1}\bigr) ,
\qquad
B = \frac{\sinh\bigl(\sqrt{p}\,\xi_{1}\bigr)}{b\sqrt{p}} ,
\qquad
C = b\sqrt{p}\,\sinh\bigl(\sqrt{p}\,\xi_{1}\bigr) , \]
avec $AD - BC = \cosh^{2} - \sinh^{2} = 1$. Pour une couche d'épaisseur $D_{\text{ép}}$,
$\xi_{1} = D_{\text{ép}}/\sqrt{a}$ et l'on retrouve l'expression classique du quadripôle homogène.

Cette limite constitue le critère de validation de toute expression quadripolaire de couche graduée.

\section{Indétermination du retour à la profondeur physique}

\subsection{Nature de l'indétermination}

La reconstruction d'un profil dans l'axe physique exige la relation entre $z$ et $\xi$. Or cette
relation dépend de la diffusivité, qui figure parmi les grandeurs recherchées. Toute tentative
d'inversion directe conduit à une définition récursive.

\subsection{Famille paramétrée}

\[ z(\xi) = \int_{0}^{\xi} \lambda^{q}\, c^{\,q-1}\, b^{\,1-2q}\, du , \]
identité valable pour toute valeur de $q$, l'exposant de $\lambda$ se réduisant à $1/2$ et celui de $c$
à $-1/2$ après regroupement.

\begin{center}
\begin{tabular}{@{}cl@{}}
\toprule
$q$ & Propriété supposée constante \\
\midrule
$0$ & capacité thermique volumique \\
$1/2$ & diffusivité \\
$1$ & conductivité \\
\bottomrule
\end{tabular}
\end{center}

Lorsque le préfacteur $\lambda^{q}c^{\,q-1}$ est constant, il sort de l'intégrale et la quadrature se
réduit à la primitive d'une puissance de $b$. En dehors de ces trois valeurs, $q$ joue le rôle d'une
hypothèse de fermeture intermédiaire sans correspondre à une propriété matérielle particulière.

\subsection{Information supplémentaire requise}

Deux fonctions indépendantes sont nécessaires à la définition complète d'un milieu thermique. La mesure
unidimensionnelle en fournit au plus une, l'effusivité transformée. Une information indépendante est
donc requise : mesure directe d'une propriété, mesure indépendante de l'épaisseur, ou hypothèse de
fermeture explicite.

Les distributions de propriétés compatibles avec un même signal peuvent s'étendre sur plusieurs ordres
de grandeur. L'article rapporte des variations de conductivité couvrant deux décades entre les valeurs
extrêmes de $q$, pour un champ mesuré identique.

\section{Potentiel constant et trichotomie}

\subsection{Les trois régimes}

L'imposition $V = \beta$ ramène la forme normale à une équation à coefficients constants,
$s'' = \beta s$, dont le signe de $\beta$ détermine la nature des solutions.

\begin{center}
\begin{tabular}{@{}lll@{}}
\toprule
Signe & Solutions pour $s(\xi)$ & Famille \\
\midrule
$\beta > 0$ & $\cosh(\xi/\xi_{c})$, $\sinh(\xi/\xi_{c})$ & hyperbolique \\
$\beta = 0$ & fonctions affines de $\xi$ & linéaire \\
$\beta < 0$ & $\cos(\xi/\xi_{c})$, $\sin(\xi/\xi_{c})$ & trigonométrique \\
\bottomrule
\end{tabular}
\end{center}

Pour $\beta$ positif, la dérivée seconde est de même signe que la fonction et la solution diverge
exponentiellement. Pour $\beta$ négatif, elle est de signe opposé et produit une oscillation. La
famille affine constitue la limite dégénérée entre les deux et non une quatrième famille indépendante.

Le profil d'effusivité s'en déduit par $b = s^{\pm 2}$, et la longueur caractéristique vaut
$\xi_{c} = |\beta|^{-1/2}$, dont le carré est un temps de diffusion propre au profil.

\subsection{Décalage spectral}

Un potentiel constant équivaut à un décalage du paramètre spectral,
\[ p \longrightarrow p + \beta , \]
les solutions du champ devenant $\cosh\bigl(\sqrt{p+\beta}\,\xi\bigr)$ et
$\sinh\bigl(\sqrt{p+\beta}\,\xi\bigr)$. La structure mathématique demeure celle du cas homogène.

\subsection{Paramètre de courbure}

Le rapport entre le temps caractéristique du profil, $\xi_{c}^{2}$, et le temps de diffusion à travers
la couche contrôle le degré de courbure du profil. Un rapport élevé correspond à un comportement proche
de l'uniformité ; un rapport faible correspond à une courbure prononcée.

\section{Points critiques relevés}

\subsection{Statut de l'analogie spectrale}

L'auteur précise que la dénomination de Schrödinger ne repose que sur une similarité de forme. Aucune
condition d'intégrabilité de carré ne s'applique à $\psi$, les solutions sont recherchées pour des
valeurs complexes arbitraires de $p$, et aucune recherche de valeurs propres ni d'états liés n'est
menée.

\subsection{Restitution des variations profondes}

L'article établit par ailleurs que l'effusivité apparente ne rejoint la valeur du substrat que pour des
fréquences normalisées inférieures à $10^{-2}$, voire $10^{-3}$. Les dépassements présentés par
certains profils n'apparaissent pas dans les courbes correspondantes. L'auteur relie explicitement
cette restitution imparfaite à la difficulté d'obtenir une inversion thermique précise en profondeur,
sans en fournir l'explication structurelle.

\subsection{Contraste entre diffusion et propagation}

En section 4, l'auteur rapporte que les profils duaux appliqués au problème électromagnétique
produisent des spectres de réflectance identiques, alors qu'en diffusion thermique deux profils duaux
distincts ne peuvent jamais donner la même réponse à toute fréquence. Le passage de l'équation de
diffusion à l'équation d'onde s'accompagne donc, sur cet axe, d'une perte d'unicité.
\end{document}